% ============================================================
%  Section: Numerical Experiments
% ============================================================

\section{Numerical Experiments}
\label{sec:experiments}

We present three sets of numerical experiments designed to interrogate the behaviour of the
Hellinger--Kantorovich (HK) gradient flow and the Newton--Hellinger flow developed in the
preceding sections.
The first experiment establishes a baseline by comparing both flows on a
well-separated Gaussian mixture, where the identifiability conditions are comfortably satisfied.
The second set of experiments examines, on the same family of mixtures, two
structural choices that arise naturally from the gradient flow formulation: whether the flow
should be stopped after a single sequential pass over the data or allowed to continue with
bootstrapped index resampling, and whether the Fisher--Rao prior regularisation term should be
included in the Hellinger component of the flow.
The third experiment stresses both flows on a significantly harder target --- a
seven-component bivariate mixture with the topology of a cat-paw --- and uses it to expose a
metastability phenomenon that the continuation mechanism is designed to resolve.
All experiments are implemented in JAX and the code is available at
\url{https://github.com/BernardoFL/RecursiveMixtures}.

% ------------------------------------------------------------
\subsection{Experiment 1: Newton--Hellinger versus HK Flow on a
             Three-Component Gaussian Mixture}
\label{subsec:exp1}

% Setup
We begin with the simplest setting in which both algorithms are well-posed.
The target mixing measure is supported on three atoms located at means
$\mu_1 = -2.0$, $\mu_2 = 0.0$, $\mu_3 = 2.0 \in \mathbb{R}$, with standard
deviations $\sigma_1 = \sigma_3 = 0.5$, $\sigma_2 = 0.8$, and weights
$w_1 = w_3 = 0.3$, $w_2 = 0.4$.
A dataset of $n = 1{,}000$ i.i.d.\ observations is drawn from the corresponding
Gaussian mixture.
Both flows are initialised with $M = 100$ particles drawn from a
Pitman--Yor process prior $\mathrm{PY}(d, \theta, G_0)$ with discount $d = 0.2$,
concentration $\theta = 10.0$, and Gaussian base measure $G_0 = \mathcal{N}(0, 9)$.
The HK flow uses kernel bandwidth $h = 0.8$, step size $\varepsilon = 0.1$,
and Sinkhorn regularisation $\lambda = 0.05$ with Wasserstein weight $\rho = 0.1$.
The Newton--Hellinger flow uses the same particle initialisation but updates weights
only, with atoms held fixed at their initial locations, via the decaying step-size
schedule $\alpha_n = (n+1)^{-0.6}$.

% Results
\begin{description}
  \item[Results.]
  \textit{[TODO: Describe the four-panel figure output.
  Panel 1: density evolution over time for each flow.
  Panel 2: final density estimate vs.\ ground truth for each flow.
  Panel 3: particle trajectories coloured by final mode assignment.
  Panel 4: heatmap of weight dynamics across steps.
  Report qualitative behaviour --- does HK recover atom locations?
  Does Newton--Hellinger concentrate weight on the correct atoms?
  Comment on convergence speed and visual alignment with the true density.]}
\end{description}

% ------------------------------------------------------------
\subsection{Experiment 2: Bootstrap Studies on the Three-Component Mixture}
\label{subsec:exp2}

The gradient flow formulation motivates two questions about runtime choices that are not
addressed by theoretical convergence results alone.
First, the flow is defined as a continuous-time object driven by a data stream; in finite
samples it is natural to ask whether stopping after a single sequential pass over the data
exhausts the useful signal, or whether continuing beyond the data by resampling indices
provides additional regularisation.
Second, the Fisher--Rao component of the HK geometry induces a prior regularisation term in
the Hellinger gradient; its effect on finite-sample behaviour --- whether it accelerates
convergence or merely adds noise at moderate $n$ --- is not self-evident from the geometry
alone.
We address these questions through two structured bootstrap studies, both using the same
three-component bivariate Gaussian mixture target from Section~\ref{subsec:exp1}, with sample
sizes $n \in \{50, 100, 200\}$.

Each experiment uses Bayesian bootstrap resampling.
For each Monte Carlo replicate, Dirichlet-style weights are drawn over data indices and a
bootstrap dataset of length $n$ is constructed by sampling training indices with replacement
according to those weights.
The flow is then run on the resulting resampled stream.
Both studies use $M = 50$ particles, Sinkhorn iterations set to 25 per optimal transport
solve, and the same Pitman--Yor prior as Experiment~1.

\subsubsection{Study A: Truncation versus Index Continuation}
\label{subsubsec:studyA}

In the truncated condition the flow processes indices $0, \ldots, n-1$ sequentially and
terminates; exactly one epoch over the bootstrap stream.
In the continuation condition the flow runs for $\lceil c \cdot n \rceil$ steps with
$c = 2.0$: the first $n$ steps consume the data in order, after which additional steps draw
indices uniformly with replacement from the same bootstrap dataset (index continuation via
\texttt{bootstrap\_after\_data=True}).
For each $n$ in the sample-size grid, both conditions are run on the same bootstrap replicate
and their final particle configurations are visualised against the true density heatmap.
Particles are rendered with area proportional to weight; no data scatter is shown so that the
figures isolate the effect of the flow itself rather than any one realisation of the data.

% Results
\begin{description}
  \item[Results.]
  \textit{[TODO: Describe Figure~\ref{fig:bootstrap_truncation} (the multi-page PDF,
  one page per $n$, two panels per page: teal = truncated, crimson = continuation).
  At small $n$ (e.g.\ $n = 50$): does continuation recover structure that truncation misses,
  or do both collapse to a small region?
  At larger $n$: does the gap between the two conditions shrink?
  Is there evidence that continuation over-smooths or introduces spurious modes?
  Summarise the trend across sample sizes.]}
\end{description}

\subsubsection{Study B: Fisher--Rao Prior Regularisation}
\label{subsubsec:studyB}

Study B fixes $n_{\text{steps}} = n$ and $\texttt{bootstrap\_after\_data} = \texttt{False}$
in both arms, isolating the contribution of the Fisher--Rao prior term.
In the regularised arm the Sinkhorn prior functional contributes to the Hellinger gradient
with weight \texttt{prior\_flow\_weight}; in the unregularised arm this term is suppressed
entirely, even if its weight parameter is positive.
Crucially, the Sinkhorn drift toward the prior at the atom level (\texttt{use\_sinkhorn}) is
\emph{not} toggled: only the weight-update prior term is varied.
Results are displayed as an $N \times 2$ grid (one row per $n$, teal = regularised,
royalblue = unregularised) using the same visualisation conventions as Study A.

% Results
\begin{description}
  \item[Results.]
  \textit{[TODO: Describe Figure~\ref{fig:bootstrap_prior} (the $N \times 2$ grid PDF).
  Does prior regularisation visibly pull particles toward the true modes at small $n$?
  At large $n$, does its effect wash out as expected from a consistent estimator?
  Are there pathological cases where regularisation hurts --- e.g.\ particles trapped
  near the prior mean rather than the data-supported modes?
  Comment on whether the separation between the two arms is monotone in $n$.]}
\end{description}

% ------------------------------------------------------------
\subsection{Experiment 3: Metastability on the Cat-Paw Distribution}
\label{subsec:exp3}

% Setup
The third experiment is designed to expose a qualitative failure mode of the plain HK flow
and to demonstrate that index continuation resolves it.
The target is a seven-component bivariate Gaussian mixture whose components are arranged in
the topology of a cat paw: four elongated toe pads positioned on a circular arc (radius
$r = 1.6$, centred at $(0, -0.55)$, spanning $30°$ to $150°$) and three rounder palm
pads at $(-0.85, -1.45)$, $(0.85, -1.45)$, and $(0.0, -0.55)$.
The toe components have anisotropic covariance $(\sigma_x, \sigma_y) = (0.18, 0.36)$,
the side palms are isotropic with $\sigma = 0.38$, and the central palm has $\sigma = 0.55$.
Mixture weights are approximately $0.11$ per toe and $0.15$ / $0.25$ for the side and
central palms respectively.
This target is substantially harder than the three-component mixture: the seven components
span a compact two-dimensional region, the toe pads are narrow and anisotropic, and
the layout creates local minima in the energy landscape that can trap particles during
the sequential data pass.

A dataset of $n = 1{,}000$ observations is drawn from the cat-paw distribution.
Three variants of the HK flow are initialised with the \emph{same} $M = 50$ particles drawn
from a Pitman--Yor prior over atom locations, so that differences across variants are
attributable solely to the algorithmic choices and not to randomness in initialisation.
The three variants are:
(a) $n$ flow steps with Fisher--Rao prior regularisation enabled;
(b) $n$ steps with prior regularisation disabled; and
(c) $n + k$ steps with prior regularisation enabled, where the first $n$ steps consume the
data sequentially and the remaining $k = 1{,}000$ steps perform index continuation via
uniform resampling from the same $n$ observations.
All variants share step size $\varepsilon = 0.01$, kernel bandwidth $h = 0.4$, and Sinkhorn
regularisation $\lambda = 0.05$.

Results are presented as a three-panel figure showing the true density heatmap alongside a
data scatter and, for each variant, the final particle configuration with marker area
proportional to weight.

% Results
\begin{description}
  \item[Results.]
  \textit{[TODO: Describe Figure~\ref{fig:paw_comparison} (the three-panel PDF:
  (a) prior on, n steps; (b) prior off, n steps; (c) prior on, n+k steps).
  For panel (a): do particles recover most of the seven components, or do some
  get stranded between modes?  Is the central palm over-represented?
  For panel (b): without prior regularisation, does the flow concentrate faster but
  miss modes, or does it spread more uniformly?  Is there evidence of mode collapse?
  For panel (c): does index continuation visibly resolve metastability ---
  do particles that were stranded in (a) migrate to the correct modes after the
  continuation phase?  How much weight redistribution is visible?
  This is the main qualitative claim of the metastability experiment; describe it carefully.]}
\end{description}
